\sscp{Labials *p, *b, *ma-p, *-mp-, *-mb-}{13-01-02}
The behaviour of labials has been suggested as being of subgrouping
relevance for CMP \citep{Blust1981,Blust1993,Blust2009}.
\trf{13-02} shows the labials in various combinations for each subgroup.
An occurrence of one does not constitute a pattern,
so while we may find dozens of reflexes of *p
and *b, confident reflexes of nasal + labial combinations for
many languages might be quite sparse.
For PMP *p the direction of sound change is assumed to generally
be *p > \it{ɸ > f > h} > {\0}.
So if there is a language in a subgroup with \it{ɸ},
that is the value given for that subgroup, and so forth.
\tsc{Ambon-Seram} is given values of \it{p-,
-f-, -f}, as the values for those positions,
even though most \tsc{Ambon-Seram} languages in most nodes
and in most words have \it{h} or {\0} in all positions (see \srf{07-03-03}).

\begin{table}
	\centering\tcp{Summary of labials in each subgroup}{13-02}
%\begin{adjustbox}{width=\textwidth}
	\st{5.75pt}\begin{tabular}{llll|ll|ll|ll}\lsptoprule
PMP				&	*{\#}p			&	*-p-		&	*p{\#}				&	*{\#}b		&	*-b-			&	*ma-p			&	*ma-b			&	*-mp-			&	*-mb-\\	\midrule
\tsc{BL}	&	{\hsh}p		&	{\hd}p		&	{\hr}{\0}			&	{\hsh}b, β&	{\hd}β, b	&	{\hr}map	&	{\hr}mab	&	{\hd}mp?	&	{\hd}mb\\
\tsc{TB}	&	{\hsh}h		&	{\hd}h		&	{\hr}p				&	{\hsh}b		&	{\hd}b		&	{\hr}map	&	{\hr}mab	&	{\hd}mb		&	{\hd}mb\\
\tsc{CT}	&	{\hsh}p		&	{\hd}p		&	{\hr}{\0}			&	{\hsh}b		&	{\hd}b		&	{\hr}mp		&	{\hr}mb		&	{\hd}mb?	&	{\hd}mb\\
\tsc{Aru}	&	{\hsh}ɸ		&	{\hd}ɸ		&	{\hr}ɸ				&	{\hsh}p		&	{\hd}p		&	{\hr}?		&	{\hr}?		&	{\hd}mb		&	{\hd}mb\\
STB				&	{\hsh}p		&	{\hd}p		&	{\hr}p				&	{\hsh}b		&	{\hd}b		&	{\hr}map	&	{\hr}mab	&	{\hd}mb		&	{\hd}mb\\
Banda			&	{\hsh}{\0}&	{\hd}{\0}	&	{\hr}p				&	{\hsh}f		&	{\hd}f		&	{\hr}mb		&	{\hr}mb		&	{\hd}mb		&	{\hd}mb\\
\tsc{AS}	&	{\hsh}p		&	{\hd}f		&	{\hr}f, {\0}	&	{\hsh}b		&	{\hd}b		&	{\hr}map	&	{\hr}maβ	&	{\hd}mp		&	{\hd}mp\\
\tsc{SB}	&	{\hsh}p		&	{\hd}p		&	{\hr}{\0}			&	{\hsh}b		&	{\hd}b		&	{\hr}mb		&	{\hr}mb		&	{\hd}mb		&	{\hd}mb\\
%\tsc{Tal}	&	{\hsh}h		&	{\hd}h		&								&	{\hsh}f 	&	{\hd}f		&	{\hr}mb		&	{\hr}mb		&	{\hd}mb		&	{\hd}mb\\
\lspbottomrule
	\end{tabular}
%\end{adjustbox}
\end{table}

\citet[26]{Blust1981} says,
“One of the more striking innovations which appears to be shared
exclusively by CMP languages is the merger of PAN *-mb-
and *-mp- as PCMP *mb.” In our data these clusters do appear
to merge in most subgroups,
but they appear to remain distinct in \tsc{Bima-Lembata}.
Note also that in \tsc{Ambon-Seram} the merger is devoiced to
**mp, even in languages that retain the voicing of *b = \it{b},
so the devoicing of *mp/*mb > **mp is significant.
The proposed merger of *mp/*mb > **mb thus cannot be accepted
at the level of putative PCMP.
While we could posit *mp > **mb to account for the reflexes in
several subgroups, this change simply consists of the voicing of a plosive after
a homorganic nasal and is so common as to be of negligible subgrouping value.

Furthermore, in many languages this change relies on only one
witness, usually reflexes of *umpu/*ampu/*əmpu ‘grandparent, grandchild’.
This issue is further complicated by forms in some languages
which cannot derive from a medial *NC cluster,
such as *upu > Buru \it{opo} ‘grandparent, grandchild (reciprocal term)’.
In Buru *p = \it{p}
and *mp > \it{b} are regular.
*umpu would derive an unattested *[obo]; compare that with PMP *ta-umpu ‘grandparent/grandchild’ (cf.
Dempwolff gloss ‘master, lord’) > Buru \it{tobo-n} ‘master, lord’.
Thus the extent to which *mp > **mb can be established in many
languages as a regular sound correspondence with any confidence is questionable.
The low frequency of forms,
the optional nasal, and exceptions like Buru diminish its reliability
and value for subgrouping.

\citet[72]{Blust2009} also appeals to what he calls the “postnasal
voicing of stops” where PMP *ma-p > \it{b, mb}.
Blust describes this as a “three-step change that involves: 1)
loss of the vowel in PMP *ma- ‘stative’,
2) postnasal voicing of a base-initial voiceless stop,
and 3) loss of the conditioning nasal.” We note that the “postnasal
voicing of stops” is part of a broader pattern involving nasal-plosive
sequences including PMP *mp/*mb, *nt/*nd, *ŋk/*ŋg.
In this section we focus on the labials,
and specifically address the claim that it is an exclusively
shared innovation in support of PCMP.
This topic was introduced briefly in \srf{01-07},
where it was noted there are actually several patterns in the
region, not just the one described by Blust.
We have identified six patterns in the Wallacean data.
A pattern such as *ma-p can:

\begin{enumerate}
	\item	Retain the labial with no voicing,
				and lose the full prefix along with the nasal,
				as in *ma-panas ‘hot’ > \it{hanas} ‘hot’ (Tetun).
				This may simply be a retention of *panas ‘hot’ without the historical
				stative *ma- prefix and thus may not actually be directly relevant to this discussion.
	\item	Retain the nasal of the prefix
				and lose the vowel
				and the labial onset of the root,
				as in *ma-panas ‘hot’ > Meto \it{manas} (‘sun’ in most Meto languages;
				‘hot’ in Kusa-Manea); *ma-putiq ‘white’ > Meto \it{mutiʔ} ‘white, silver’.
	\item	*ma-p > \it{p},
				but not merge with *p,
				and no post-nasal voicing,
				as in Dhao *ma-panas > \it{pana} ‘hot’,
				*ma-putiq > \it{puɖ͡ʐi} ‘white,
				silver’, *ma-pənuq > \it{pənu} ‘full’; (Dhao *p > {\0},
				\it{p} as in *pajay > \it{are} ‘rice’,
				*pusuq > \it{usu} ‘heart’,
				*pitu > \it{piɖ͡ʐu} ‘seven’).
				Or Alune *ma-putiq > \it{puti} ‘white’,
				*ma-pənuq > \it{penu} ‘full’; (Alune *p > {\0} as in *pajay >
				\it{ala} ‘rice’, *pija > \it{ila} ‘how many’,
				*punti > \it{uri} ‘banana’,
				*hapuy > \it{au-e} ‘fire’; note that Alune retains voicing in
				*b = \it{b}; see \srf{10-02}).
	\item	*ma-p > \it{b} becoming voiced (= “post-nasal voicing of
				stops” as described by Blust),
				as in Buru *ma-panas ‘hot’ > \it{bana} ‘fire’,
				*ma-putiq > \it{boti} ‘white,
				ripe (grain)’, *ma-pənuq ‘full’ > \it{ben{\tl}beno} ‘sag from weight
				(of heavily laden fruit)’.
	\item	Retain the nasal,
				retain the labial, lose the vowel,
				become voiced as in Kambera *ma-panas > \it{mbana},
				‘hot, warm’, \it{mbanah-u} ‘hot’.
	\item	Retain both the full historical prefix with the vowel
				and the labial, and stay voiceless as in *ma-panas ‘hot’ > Gorom
				\it{mahanas} ‘hot’, Geser \it{mfanas} ‘hot’, Yamdena \it{mefanas} ‘hot’; *ma-putiq
				‘white’ > Gorom \it{mahuti} ‘white’,
				Geser \it{mafuti} ‘white’, Yamdena \it{mafuti} ‘white’.
\end{enumerate}

It is possible to make a further higher-level generalisation
that numbers 3-5 in the list above all share the process of fusion
that results in a full merger.\footnote{
		We thank a reviewer for clarifying that these three patterns all share the same process of fusion.}
We have laid them out as separate above as this list highlights
clearly the different surface forms that people (including descriptive
linguists and anthropologists) will encounter.

We have no data in which *ma-p > \it{mab} with voicing of the labial.
However *ma-b = \it{mab} in *ma-bəRat ‘heavy’ > Yamdena \it{maberat}
‘heavy’, is not directly relevant to the issue in focus,
other than showing different reflexes of *ma-p
and *ma-b in the same language (compare with examples in the previous paragraph).
The data in \trf{13-02} show:

\begin{itemize}
	\item[a.]	that the distinction between *ma-p
						and *ma-b must be reconstructed for the region;
	\item[b.]	they are often different than the reflexes for *-mp-
						and *-mb-;
	\item[c.]	the full sequences must be reconstructed to
						account for all the data;
	\item[d.]	neither the proposed merger of *mp/*mb > **mb nor the “postnasal
						voicing of stops” can be accepted as having subgrouping value
						for all the Wallacean subgroups in the target region.
						They do not provide support for the putative PCMP.
\end{itemize}

We see shared retentions (not relevant for subgrouping),
but no splits, mergers or innovations that are shared by all
subgroups, other than the merger of
*mp/*mb/*ma-p/*ma-b > **mb for both \tsc{Sula-Buru} and Banda.\footnote{
		In \tsc{Sula-Buru} languages, the surface manifestation of **mb > \it{b},
		whereas Banda retains the nasal-consonant cluster **mb = \it{mb}.}
However, given that those share very little
else in their historical phonologies,
we see this merger as parallel development,
not a shared innovation.

The situation is actually more complicated
and more interesting than just the patterns noted above.
For example, closely related languages
and even single languages can have more than one pattern for
the same historical sequence,
so forms commonly reconstructed with the stative prefix *ma-,
should also be investigated without such a prefix.
Such complexity is shown for a number of closely related Meto
languages in example \qf{13-01}.
Investigators in the region need to be careful to not “find one and be done”.

\noindent{\wg\st{6pt}\begin{tabularx}
{\linewidth}{>{\hspace{-\tabcolsep}}l<{\hspace{-\tabcolsep}}
>{\hspace{-\tabcolsep}\enspace}ll
>{\raggedright\arraybackslash}X}
\lsptoprule
%\tbx{13-1}	&	PMP	&	*(\tbr{ma})-\tbr{p}anas	&	hot	\\	\midrule
%\endfirsthead
\tbx{13-01}	&	PMP	&	*(\tbr{ma})-\tbr{p}anas	&	hot	\\	\midrule
%\endhead
	&	Meto (most documented varieties)	&	\tbr{m}anas	&	sun	\\
	&	Kusa-Manea	&	\tbr{m}anas	&	hot	\\
	&	Timaus	&	na-\tbr{m}ana	&	it is midday	\\
	&	Meto (most documented varieties)	&	na-\tbr{h}ana	&	cook (< *panas)	\\
	&	Kotos Amarasi	&	k-\tbr{h}anas	&	drought, hot/dry season (*ka-panas)	\\
\lspbottomrule
\end{tabularx}\mbox{}\\}

\tsc{Southwest Maluku} languages of the \tsc{Kisar-Luang} node of \tsc{Timor-Babar}
(\srf{04-06-02}) show variety both for the same etymon,
and also different sound correspondences for different etyma,
as illustrated in examples \qf{13-02} and \qf{13-03}.
Note that these languages normally have *p > {\0} (see \trf{04-01})
indicating that reflexes with \it{p} are from intermediate **mp.

\noindent{\wg\st{6pt}\tblr{>{\hspace{-\tabcolsep}}l<{\hspace{-\tabcolsep}}
>{\hspace{-\tabcolsep}\enspace}
lll}
\lsptoprule
%\tbx{13-2}	&	PMP	&	*\tbr{ma})-\tbr{p}anas	&	hot	\\	\midrule
%\endfirsthead
\tbx{13-02}	&	PMP	&	*\tbr{ma}-\tbr{p}anas	&	hot	\\	\midrule
%\endhead
	&	Kisar	&	\tbr{m}anha	&	hot	\\
	&	Roma	&	\tbr{m}ahan	&	hot	\\
	&	Luang	&	\tbr{p}anha	&	hot (**mpanas; cf. *umpu > \it{upə} grandchild)	\\
	&	Wetan	&	\tbr{p}aana	&	hot (**mpanas)	\\	\midrule
\tbx{13-03}	&	PMP	&	*\tbr{ma}-\tbr{p}ənuq	&	full	\\
	&	Kisar	&	\tbr{p}enu	&	full	\\
	&	Roma	&	\tbr{p}empwenu	&	full (/penu{\tl}penu/)	\\
	&	Luang	&	\tbr{p}en-penu	&	full	\\
	&	Wetan	&	\tbr{p}eni	&	full	\\
\lspbottomrule
\end{tabular}\mbox{}\\}

While \citet[72]{Blust2009} sees “post-nasal voicing of stops”
as “an innovation, and one that is highly distinctive”
and “not known to be shared with any language outside eastern
Indonesia”, he nevertheless concedes that “this change,
like glide truncation, was a product of diffusion,
since some languages that show no evidence of it are interspersed
with those that do.” We agree.
(Glide truncation is discussed in the next section \srf{13-01-03}.)

Superficially it appears that \tsc{Ambon-Seram},
Banda, \tsc{Aru}, \tsc{Seram-Tanimbar-Bomberai}, and \tsc{Timor-Babar}
all share a weakening of *p > \it{ɸ, f, h}, {\0}.
While such weakening is present as an areal tendency for some
subgroups, nevertheless at the level of whole subgroups,
the picture is deceptive.
For example, \tsc{Ambon-Seram} retains initial *p = \it{{\#}p} in
some languages and some lexical items that cannot be accounted for by a historical
*ma-, nor by consonant alternations in the onset of the root from subject
indexing on verbs (e.g.
Tehoru \it{puse-n} ‘navel’, Alune \it{benu} ‘sea turtle’,
and reflexes of causative *pa- = \it{pa-} in several languages).
Banda retains final *p in *qatəp > \it{atop} ‘thatch’.
In \tsc{Seram-Tanimbar-Bomberai}, Kur retains *p = \it{p} in all positions,
and Yamdena retains final *p = \it{p{\#}} (e.g.
\it{malip} ‘laugh’, \it{-morip} ‘live’).
\tsc{Timor-Babar} retains final *p = \it{p{\#}} in SW Tanimbar.

The weakening of *p to a voiceless fricative (\it{ɸ,
f, h}) is a common sound change cross linguistically,
including in Austronesian languages in the Philippines
and in western Indonesia,
so parallel development cannot be ruled out.
Positing a higher level grouping reconciling the diversity of
the exceptions noted above would be awkward.
The exceptions word initially in some \tsc{Ambon-Seram} languages,
in all positions in Kur,
and word finally in Yamdena
and SW Tanimbar cannot be accounted for by a change in a single node.
The same change in SHWNG shows it is an areal phenomenon that
probably spread by diffusion.
Even appealing to a linkage would be difficult,
given that retentions of *p = \it{p} are spread unevenly all over the region.

At best, such a change would provide only very weak subgrouping evidence.
We are hesitant to propose a large subgroup on the basis of this
one common but uneven sound change alone.
Unless further shared innovations in the phonology,
morphology, and/or lexicon can be identified between these languages,
we prefer not to place them in a single subgroup.
If further evidence favours such a merger,
the resulting subgroup could be a candidate for some sort of
modified “Central Malayo-Polynesian” subgroup,
although how it links up to PMP,
and how it relates to \tsc{Bima-Lembata},
\tsc{Central Timor}, \tsc{Sula-Buru}, SHWNG and Oceanic remains elusive.

